% 326_Matzke_FFGFT_Vergleich_De_ch.tex
% Dok. 326, FFGFT-Korpus, August 2026 — Kapitel-Datei
% Wird von Sources/wr_standalone_A4/326_Matzke_FFGFT_Vergleich_De.tex eingebunden.

% --- Dokumentlokale Makros (kollisionssicher) ---
\providecommand{\K}{}\renewcommand{\K}{\textbf{[K]}}
\providecommand{\B}{}\renewcommand{\B}{\textbf{[B]}}
\renewcommand{\S}{\textbf{[S]}}
\providecommand{\Q}{}\renewcommand{\Q}{\textbf{[Q]}}
\providecommand{\X}{}\renewcommand{\X}{\textbf{[X]}}
\providecommand{\SETZ}{}\renewcommand{\SETZ}{\textbf{[SETZUNG]}}
\providecommand{\lP}{}\renewcommand{\lP}{\ell_{\mathrm{P}}}
\providecommand{\mP}{}\renewcommand{\mP}{m_{\mathrm{P}}}
\providecommand{\TP}{}\renewcommand{\TP}{T_{\mathrm{P}}}
\renewcommand{\TH}{T_{\mathrm{H}}}
\providecommand{\rs}{}\renewcommand{\rs}{r_{\mathrm{s}}}
\providecommand{\SBH}{}\renewcommand{\SBH}{S_{\mathrm{BH}}}
\providecommand{\xig}{}\renewcommand{\xig}{\xi}
\providecommand{\Tmass}{}\renewcommand{\Tmass}{\tilde{T}}
\providecommand{\Ebit}{}\renewcommand{\Ebit}{E_{\mathrm{bit}}}
% --- Farbboxen ---
\definecolor{ffgftblue}{RGB}{30,90,160}
\definecolor{matzkegreen}{RGB}{0,120,60}
\definecolor{converge}{RGB}{160,60,0}
\definecolor{diverge}{RGB}{140,0,40}
\definecolor{warn}{RGB}{100,0,120}
\tcbset{
  base/.style={breakable,enhanced,boxrule=0.6pt,arc=3pt,
    fonttitle=\bfseries\small,coltitle=white,top=4pt,bottom=4pt,left=6pt,right=6pt},
  ffgftbox/.style={base,colback=blue!5,colframe=ffgftblue,colbacktitle=ffgftblue},
  matzbox/.style={base,colback=green!5,colframe=matzkegreen,colbacktitle=matzkegreen},
  convbox/.style={base,colback=orange!8,colframe=converge,colbacktitle=converge},
  divbox/.style={base,colback=red!6,colframe=diverge,colbacktitle=diverge},
  warnbox/.style={base,colback=violet!6,colframe=warn,colbacktitle=warn}
}

\begin{tcolorbox}[ffgftbox, title={Zeichenerklärung}]
  \K\ Kernresultat\quad \B\ algebraisch bewiesen\quad \S\ Skizze\quad \Q\ Quelle\quad \X\ ausgeschlossen\quad \SETZ\ Axiom
\end{tcolorbox}

% ============================================================
\section*{Zusammenfassung}

Douglas Matzke leitet im Beitrag \enquote{Black Holes from Hyperbits} (IPI 2026)
die gesamte Schwarze-Loch-Physik aus dem Clifford-Algebra-Axiom $e_i^2=+1$ ab,
ohne Allgemeine Relativitätstheorie und ohne freie Parameter.
Zentral dabei: die Masse eines Bits folgt aus Landauers Prinzip
($E_{\mathrm{bit}} = k_B T_P \ln 2$) plus $E=mc^2$ und ist damit eine
\emph{universelle} Zahl — $m_{\mathrm{bit}} = m_P \ln 2/(2\pi) \approx 0{,}11\,m_P$,
unabhängig vom physikalischen System.

FFGFT (Fundamentale Fraktale Geometrische Feldtheorie) kommt zu einer fundamental
anderen Einsicht: Bitwerte sind \emph{systemabhängig}.
Die minimale Informationseinheit ist das Windungsquant $\Delta w = 1$ auf $T^4$
(Dok.\ 257), und ihre Energie ist $E_{\mathrm{bit}} = \hbar c/L$ — abhängig von der
physikalischen Skala $L$ des Systems.
Landauers $k_B T \ln 2$ ist der Spezialfall einer bestimmten Skala und Temperatur,
kein universelles Gesetz (A271, A272, A273, Dok.\ 302).

Dieses Dokument vergleicht beide Frameworks und arbeitet die Differenz in der
Landauer-Interpretation als theoretisch wesentlich heraus.



% ============================================================
\section{Warum Matzkes Algebra so einfach wirkt}
% ============================================================

Bevor die technischen Vergleiche folgen, lohnt eine narrative Einordnung:
Woher kommt der ästhetische Reiz des Hyperbit-Zugangs -- und wo ist er
erkauft? Beides gehört zusammen.

\subsection*{Der Ausgangspunkt: ein Vorzeichen}

Matzke beginnt mit $e_i^2 = +1$. Das ist ein \emph{Bit} -- etwas, das
quadriert Eins ergibt, also zwei Zustände hat. Kein Raum, keine Zeit,
keine Metrik. Nur die Feststellung: es gibt unterscheidbare Dinge.

Der eigentliche Trick ist, dass die Clifford-Algebra darauf sofort zu
arbeiten beginnt. Nimmt man mehrere solche $e_i$ und lässt sie
antikommutieren, entsteht aus dem Nichts eine Struktur mit Dimension,
Orientierung und Drehungen. Hineingesteckt wurde nichts als
\enquote{es gibt Bits, und sie sind voneinander unabhängig} --
heraus kommt Geometrie. Das ist der erste Grund für die Schönheit:
\textbf{Geometrie als Buchhaltung über Unterscheidbarkeit},
nicht als Vorannahme.

\subsection*{Warum ausgerechnet $\mathrm{Cl}(6)$}

Der zweite schöne Schritt: $\mathrm{Cl}(6)$ zerfällt in genau drei
Witt-Paare -- drei Päckchen aus je einem Erzeuger und einem Vernichter.
Nach Furey ist das genau die Struktur einer Fermion-Generation:
drei Farben, ein Lepton. Die Drei ist nicht gewählt, sie \emph{fällt
heraus}. Wer je versucht hat zu erklären, warum es drei Farben gibt,
weiß, wie befriedigend das ist.

\subsection*{Der Schwarzloch-Teil}

Nun der eigentliche Coup. Bekenstein sagt: Ein Schwarzes Loch hat
Entropie proportional zur \emph{Fläche}, ein Bit pro vier
Planck-Flächen. Matzke dreht das um: Wenn ein Bit eine Fläche ist,
dann ist eine Ansammlung von Bits ein Volumen -- und die Frage
\enquote{ab wann kollabiert das?} wird zu einer Frage über Zahlen,
nicht über Feldgleichungen. Er setzt Compton-Wellenlänge gleich
Schwarzschild-Radius (der klassische Punkt, wo Quantenmechanik und
Gravitation sich treffen) und bekommt den Schwarzschild-Radius heraus --
\textbf{ohne die Einstein-Gleichungen je zu benutzen}.
Für $20\,M_\odot$: $59{,}1$~km, exakt richtig.

Und dann die Zahl, die den Reiz ausmacht: Die Stabilitätsschwelle liegt
bei $6{,}41$~Bit. $\mathrm{Cl}(6)$ hat $6{,}000$~Bit. Das Universum
besteht aus Fermionen, die exakt $0{,}41$~Bit unterhalb der
Kollapsschwelle sitzen. Materie ist stabil, weil die kleinste mögliche
Fermion-Algebra gerade eben nicht kollabiert. Das ist eine schöne
Geschichte -- knapp, dramatisch, mit einer konkreten Zahl.

\subsection*{Wo die Eleganz erkauft ist}

Diese $6{,}41$ hängt an einer Setzung. Um von Bits zu Massen zu kommen,
braucht Matzke einen Umrechnungskurs -- was wiegt ein Bit? Er nimmt
Landauer ($E = k_BT\ln 2$) und wertet ihn bei Planck-Temperatur aus.
Heraus kommt $m_{\mathrm{bit}} = m_P\ln 2/(2\pi)$, eine \emph{universelle}
Bit-Masse.

Nur: Landauer ist eine Aussage über thermische Träger-Ensembles, nicht
über abstrakte Information (A271--A273). Und die Planck-Temperatur ist
nicht die Temperatur irgendeines physikalischen Systems -- sie ist eine
Energieskala. Die Wahl $T = T_P$ ist eine Setzung, die sich nicht selbst
begründet. Die Zahl $6{,}41$ ist damit exakt so universell wie diese
Setzung -- und Dok.~325 zeigt, dass die Hawking-Temperatur ganz ohne sie
auskommt.

\begin{tcolorbox}[convbox, title={Die kürzeste Fassung}]
  Matzkes Algebra ist schön, weil sie zeigt, wie weit man mit reiner
  Kombinatorik kommt: Aus \enquote{es gibt Bits} folgen Dimension,
  drei Farben, eine Fermion-Generation, der Schwarzschild-Radius.
  Das ist echte Reduktion, keine Umbenennung.\\[4pt]
  Sie ist \emph{einfach}, weil sie an drei Stellen aufhört zu fragen:
  Bekenstein, Compton$\,=\,$Schwarzschild und Landauer bei $T_P$ werden
  als Brücken gesetzt, nicht hergeleitet. Drei geliehene Brücken zu einer
  Struktur, die intern sehr sparsam ist.\\[4pt]
  Der Kontrast zur FFGFT liegt genau dort: Die FFGFT zahlt den Preis
  vorne ($\xi$ herleiten, $L_0$, den Torus rechtfertigen) und braucht
  danach keine Brücken -- dafür ist der Einstieg schwerer. Matzke zahlt
  vorne fast nichts und leiht sich später drei Mal. Beide Rechnungen
  gehen auf; sie sehen nur an verschiedenen Stellen elegant aus.
\end{tcolorbox}

% ============================================================
\section{Ausgangspunkte}

\begin{tcolorbox}[matzbox, title={Matzke}]
  Einziges Axiom: $e_i^2 = +1$ (Hyperbit). \SETZ\\
  Physikalische Einheiten über drei externe Brücken:
  Bekenstein ($4\lP^2$ pro Bit), Compton–Schwarzschild ($\lP$ als Längenskala),
  Landauer–Einstein ($m_{\mathrm{bit}} = m_P \ln 2/(2\pi)$).
\end{tcolorbox}

\begin{tcolorbox}[ffgftbox, title={FFGFT}]
  Einziger Parameter: $\xi = 4/30000$, ableitbar als kleinster Eigenwert des
  Spektraloperators $\hat{F}$ auf $T^4/\mathbb{Z}_3$ \K\ (Dok.\ 322).
  Aus $\xi$ folgen $G$, $\lP$, $L_0 = \xi\cdot\lP$, $\alpha^{-1} = 137{,}036$,
  alle Leptonenmassen — ohne Brücken-Axiome \K\ (Dok.\ 180, 011, 006).\\
  Minimale Informationseinheit: Windungsquant $\Delta w = 1$, Energie $E_{\mathrm{bit}} = \hbar c/L$
  (\emph{skalenabhängig}) (Dok.\ 257, A271).
\end{tcolorbox}

% ============================================================
\section{Die Landauer-Frage: universelle vs.\ systemabhängige Bitwerte}

Dies ist der theoretisch tiefste Unterschied zwischen beiden Frameworks.

\subsection{Matzkes Position: universeller Bitwert}

Matzke setzt die Landauer-Energieskala bei der Planck-Temperatur $T_P$ an:
\begin{equation}
  E_{\mathrm{bit}} = k_B\,T_P\,\ln 2 \quad\Rightarrow\quad
  m_{\mathrm{bit}} = \frac{m_P\,\ln 2}{2\pi} \approx 0{,}110\,m_P
  \tag{Matzke, Kap.\ 6}
\end{equation}
Dieser Wert ist \emph{universell}: jedes Bit trägt dieselbe Masse,
unabhängig vom System (Elektron, Schwarzes Loch, Vakuum).
Daraus folgt die Stabilitätsschwelle $n_{\mathrm{thresh}} = 2\pi/(\sqrt{2}\ln 2) = 6{,}41$
und die Schwarzschild-Relation direkt.

\subsection{FFGFT-Position: systemabhängige Bitwerte}

FFGFT versteht Landauers Prinzip als Aussage über \emph{Phasenraumgebiete} eines
thermischen Träger-Ensembles, nicht über abstrakte Information (A271):

\begin{quote}
  \itshape
  Gezählt wird nach Unterscheidbarkeit (Orthogonalität), nicht nach Energie.
  Die Thermodynamik zählt Gebiete und liest sie nicht — es gibt keinen
  thermodynamischen Wahrheitsdetektor.
  Hinreichend für die Wärmeproduktion ist die \textbf{Nicht-Bijektivität der
  Gebietsabbildung}, nicht die logische Irreversibilität.
\end{quote}

\medskip
Die fundamentale Bit-Energie in FFGFT ist (Dok.\ 257, Dok.\ 302):
\begin{equation}
  E_{\mathrm{bit}} = \frac{\hbar c}{L}
  \tag{FFGFT, Dok.\ 257}
\end{equation}
wobei $L$ die physikalische Skala des Systems ist.
Landauers $k_B T\ln 2$ ist der Grenzfall bei einer bestimmten Skala und Temperatur
(A273: \enquote{zwei Böden — der thermische $k_BT$ und der quantengeometrische
$\hbar c/L$; maßgeblich ist der größere}).

Die Bit-Energie variiert damit zwischen den Skalen:

\begin{center}
\small
\begin{tabular}{@{}lll@{}}
  \toprule
  System & Skala $L$ & $E_{\mathrm{bit}} = \hbar c/L$ \\
  \midrule
  Planck-Skala & $\lP = 1{,}62\times10^{-35}$ m & $m_P c^2 = 1{,}22\times10^{19}$ GeV \\
  FFGFT-Boden & $L_0 = \xi\cdot\lP = 2{,}15\times10^{-38}$ m & $E_0/\xi \approx 56$ PeV \\
  Charakteristisch & $L_e = \hbar c/E_0 \approx 26{,}85$ fm & $E_0 = 7{,}35$ MeV \\
  Kernphysik & $\sim 1$ fm & $\sim 197$ MeV \\
  Raumtemperatur-Bit & $L_T \sim k_BT/(\hbar c) \sim\mu$m & $\sim 3\times10^{-21}$ J = $18\,\mathrm{meV}$ \\
  \bottomrule
\end{tabular}
\end{center}

\begin{tcolorbox}[warnbox, title={FFGFT-Kritik an Matzkes Landauer-Brücke}]
  Matzke setzt $T = T_P$ in $E_{\mathrm{bit}} = k_BT\ln 2$.
  Das liefert einen \emph{universellen} Bitwert.
  FFGFT widerspricht: Landauers Schranke gilt für ein thermisches Gleichgewichts-Ensemble
  bei Temperatur $T$ des Bades — dies ist eine Aussage über den \emph{Träger},
  nicht über abstrakte Information (A272).
  Die Planck-Temperatur $T_P$ ist nicht die Temperatur des physikalischen Systems
  (Elektron, Proton, Schwarzes Loch im Kosmos); sie ist nur eine Energie-Skala.
  Landauer gilt wegen der Mechanik (Liouville), nicht wegen einer privilegierten Skala.\\[4pt]
  \textbf{Konsequenz}: Matzkes universelle Stabilitätsschwelle
  $n_{\mathrm{thresh}} = 6{,}41$ hängt wesentlich an der Setzung $T = T_P$.
  Bei systemabhängiger Bit-Energie wäre $n_{\mathrm{thresh}}$ selbst systemabhängig. \S
\end{tcolorbox}

\subsection{Zwei Böden, ein Prinzip}

FFGFT unterscheidet (A273, Dok.\ 290):

\begin{description}
  \item[Statisch] Bit $\equiv$ Windungsquant $\Delta w = 1$ auf $T^4$:
    skalenuniversell, energieunabhängig als Zählgröße;
    Energie $E_{\mathrm{bit}} = \hbar c/L$ ist skalenspezifisch.
  \item[Thermodynamisch] Irreversibler Löschvorgang:
    $Q \geq k_B T\ln 2$ — Kosten in Joule, abhängig von der Badtemperatur $T$.
    Dieser Preis hat mit der informationstheoretischen Menge nichts zu tun
    (A272: \enquote{einem Informationsbit kann keine Energie zugeordnet werden,
    nur seinem Träger}).
  \item[Kinematisch] Reversible Bit-Flip-Zeit:
    $\tilde{T} = \hbar/E_{\mathrm{bit}}$;
    Margolus-Levitin $\tau \geq \pi\hbar/(2E)$;
    Energie setzt die Uhr ohne Verlust (Dok.\ 290).
\end{description}

Die vier Kleidungen einer Bit-Energie bei der charakteristischen Skala $E_0 = 7{,}35$ MeV:
\begin{equation}
  m_0 c^2 = E_0 = k_B T_0 = \hbar/\tilde{T}
  \quad\Rightarrow\quad
  T_0 = E_0/k_B \approx 8{,}5\times10^{10}\,\mathrm{K}
\end{equation}
Landauer bei Raumtemperatur liefert $Q \approx 3\times10^{-21}$ J $\approx 18$ meV —
acht Größenordnungen kleiner als $E_0$.
Die Aussagen \enquote{was das Bit ist} ($E_0 = 7{,}35$ MeV)
und \enquote{was das Löschen kostet} ($k_BT\ln 2 = 18$ meV bei 300 K)
gehören verschiedenen Schichten an (Dok.\ 290, A271).

% ============================================================
\section{Bekenstein-Entropie und Planck-Voxel}

\begin{tcolorbox}[matzbox, title={Matzke}]
  $1\,\mathrm{Bit} = \mathrm{Cl}(6) = 3$ Witt-Paare $= 1$ Planck-Voxel $4\lP^3$. \B\\
  $\SBH = A/(4\lP^2) = V_{\mathrm{shell}}/(4\lP^3)$;
  Bekbits als Grundeinheit, $\Omega = 2^b$.
\end{tcolorbox}

\begin{tcolorbox}[ffgftbox, title={FFGFT}]
  Minimale Informationseinheit: Windungsquant $\Delta w = 1$ auf $\pi_1(T^4) \cong \mathbb{Z}^4$.
  Diskretheit topologisch (keine halbe Windung), nicht aus einer Kachelung (Dok.\ 302).\\
  Planck-Voxel existieren implizit als Orbifold-Fixpunkte.
  Bekenstein: $\SBH = A/(4\lP^2)$ als Konsequenz der Bekenstein-Schranke auf der
  Saturierungsfläche \B.\\
  Holographische Richtung: $N$-dimensional $\to$ 3D (wie Matzke), nicht $3D \to 2D$.
\end{tcolorbox}

\begin{tcolorbox}[convbox, title={Konvergenz: Fermion-Generation = algebraische Grundeinheit}]
  Matzke: $\mathrm{Cl}(6) = 3$ Witt-Paare $=$ Planck-Voxel $=$ eine Fermion-Generation (Furey). \B\\
  FFGFT: $T^4/\mathbb{Z}_3$-Orbifold, drei $\mathbb{Z}_3$-Sektoren $\to$ $\mathrm{SU}(3)_c$
  $\to N_c = 3$ algebraisch notwendig $\to$ eine Fermion-Generation. \B\ (Dok.\ 321)\\
  Beide finden: drei-fache Struktur, algebraisch erzwungen, kein freier Parameter.
\end{tcolorbox}

\begin{tcolorbox}[divbox, title={Divergenz: Tiefe der Ableitung}]
  Matzke setzt $\mathrm{Cl}(6)$ als physikalische Einheit voraus (\SETZ\ über Furey-Identifikation).\\
  FFGFT leitet die Dreier-Struktur aus $\mathbb{Z}_3$-Trialisierung ab:
  $N_c = 3$ folgt algebraisch notwendig \B\ (Dok.\ 321, R79–R80).
\end{tcolorbox}

% ============================================================
\section{Schwarzschild-Radius und Bekenstein-Entropie}

Matzkes Novel Result \#1:
\begin{equation}
  r_s = \frac{2}{\sqrt{\pi}}\,\sqrt{\SBH}\,\lP
  \tag{Matzke}
\end{equation}
Der Faktor $2/\sqrt{\pi}$ kommt vom Übergang Hyperwürfel $\to$ Kugel
(\enquote{rounding the square}).

FFGFT:
\begin{align}
  G &= \frac{\xi^2}{4m_e},\quad \lP = \sqrt{\frac{\hbar G}{c^3}},\quad
  r_s = \frac{2GM}{c^2}
  \tag{FFGFT, Dok.\ 180}
\end{align}
Bekenstein: Konsequenz der Schranke $S \leq 2\pi RE/(\hbar c)$ auf der Saturierungsfläche \B.

\textbf{Numerisch}: 20\,$M_\odot$ $\Rightarrow$ $r_s = 59{,}1$ km (beide exakt).

\begin{tcolorbox}[convbox, title={Konvergenz}]
  Beide leiten $r_s$ ohne ART-Feldgleichungen ab.
  Beide zeigen: $\lP$ wird durch Selbstkonsistenz erzwungen, nicht axiomatisch gesetzt.
  Numerisch identisch.
\end{tcolorbox}

\begin{tcolorbox}[divbox, title={Divergenz}]
  Matzkes $2/\sqrt{\pi}$ folgt aus der Hypertetraeder-Geometrie.
  In FFGFT erscheint kein solcher Formfaktor; $r_s = 2GM/c^2$ ist die Standardrelation,
  und $G$ kommt aus $\xi$.
\end{tcolorbox}

% ============================================================
\section{Stabilitätsschwelle}

\begin{tcolorbox}[matzbox, title={Matzke}]
  $n_{\mathrm{thresh}} = 2\pi/(\sqrt{2}\ln 2) = 6{,}41$ Bits. \B\\
  $\mathrm{Cl}(6)$-Fermion: 6{,}000 Bits, 0{,}41 Bits unter Schwelle.\\
  Stabiles Universum, weil minimale Fermion-Algebra algebraisch unter Kollaps-Schwelle liegt.\\[2pt]
  \textbf{Abhängigkeit}: $n_{\mathrm{thresh}}$ hängt wesentlich an $m_{\mathrm{bit}} = m_P\ln 2/(2\pi)$,
  also an der Landauer-Brücke mit $T = T_P$.
\end{tcolorbox}

\begin{tcolorbox}[ffgftbox, title={FFGFT}]
  Stabile Materie folgt aus $\tilde{T}\cdot m = 1$:
  Leptonenmassen aus rationalen Invarianten $(r_i, p_i)$ des Torusraums \K\ (Dok.\ 006, 046, 320).\\
  Keine direkte Bit-Zähl-Schwelle; Stabilität aus der Orbifold-Fixpunktstruktur.\\
  $n_{\mathrm{thresh}}$ als Bit-Zahl kennt FFGFT nicht — und sollte systemabhängig sein,
  wenn Bitwerte systemabhängig sind.
\end{tcolorbox}

\begin{tcolorbox}[warnbox, title={Offene Brücke — systemabhängige Schwelle?}]
  Matzkes $n_{\mathrm{thresh}} = 6{,}41$ ist elegant, hängt aber an der Setzung
  $m_{\mathrm{bit}} = m_P\ln 2/(2\pi)$ (universeller Bitwert, Planck-Temperatur).
  Aus FFGFT-Sicht wäre die Schwelle systemabhängig:
  \begin{equation}
    n_{\mathrm{thresh}}(L) = \frac{2\pi}{\sqrt{2}} \cdot \frac{m_P c^2}{E_{\mathrm{bit}}(L)\cdot\sqrt{2}}
    = \frac{2\pi}{\sqrt{2}} \cdot \frac{m_P \cdot L}{\sqrt{2}\,\hbar/c}
  \end{equation}
  Für $L = \lP$ (Planck): ergibt dieselbe Zahl wie Matzke.
  Für $L = L_0 = \xi\cdot\lP$ (FFGFT-Boden): wäre $n_{\mathrm{thresh}}$ um Faktor $\xi$ kleiner.\\
  \textbf{Frage}: Ist $n_{\mathrm{thresh}} = 6{,}41$ die richtige Zahl für das Universum,
  oder ist sie eine Konsequenz der Wahl $L = \lP$? \S
\end{tcolorbox}

% ============================================================
\section{Hawking-Strahlung}

\begin{tcolorbox}[matzbox, title={Matzke}]
  Hawking-Quanten sind ausschließlich masselose Vakuumzustände $\{V, V_s\}$ von $\mathrm{Cl}(6)$.
  Annihilation: $P_k^+ \cdot P_-^{\mathrm{vac}} = 0$ (exakt in $\mathbb{Z}_3$). \B\\
  Information erhalten: jedes $V$ kodiert Orientierung des vernichteten Bekenstein-Voxels.\\
  $T_H = T_P/(4\,n_{\mathrm{mass}}\ln 2)$. \B\\
  Keine Bildungs-Strahlung (Bildung = Zeitumkehr der Annihilation).
\end{tcolorbox}

\begin{tcolorbox}[ffgftbox, title={FFGFT (Dok.\ 325)}]
  Vollständiger Mechanismus aus FFGFT-eigenen Bausteinen:\\
  \textbf{Temperatur}: KMS-Regel $T_H = \hbar/(k_B\tau)$ mit
  $\tau = 8\pi GM/c^3$ aus $\tilde{T}\cdot m = 1$ \K\ (Dok.\ 313 G, 325).\\
  \textbf{Quant}: $k_B T_H = \hbar c/(4\pi r_s)$ exakt — das Hawking-Quant
  ist das systemabhängige Bit der Horizontskala \K\ (Dok.\ 325).\\
  \textbf{Selektion}: nur $T_R = 0$ entkommt; Annihilation
  $P_jP_k = 0$; Confinement $=$ Hawking-Selektion \B\ (Dok.\ 321, 325).\\
  \textbf{Träger}: masselose $n_4 = 0$-Torusmode; massive Moden
  fallen zurück (Uhrenstauchung) \K\ (Dok.\ 325).\\
  \textbf{Information}: Sektorpaar $(k,-k)$ orthogonal lesbar,
  Unitarität des Spektraloperators \K\ (Dok.\ 322, 325).
\end{tcolorbox}

\begin{tcolorbox}[convbox, title={Konvergenz — zwei vollständige Mechanismen}]
  Beide Frameworks besitzen einen algebraisch vollständigen Mechanismus.
  Die Entsprechungen sind strukturell exakt (Dok.\ 325):
  Matzkes $P_k^+\cdot V_s = 0$ $\leftrightarrow$ FFGFTs $P_jP_k = 0$;
  Matzkes masseloses $V$ $\leftrightarrow$ FFGFTs $n_4=0$-Torusmode;
  Matzkes Voxel-Orientierung $\leftrightarrow$ FFGFTs Sektorpaar $(k,-k)$.
\end{tcolorbox}

\begin{tcolorbox}[divbox, title={Divergenz — die Bitwert-Grundlegung}]
  Matzkes Mechanismus ruht auf dem universellen Bitwert
  $m_{\mathrm{bit}} = m_P\ln 2/(2\pi)$; die FFGFT-Route (KMS aus
  $\tilde{T}\cdot m=1$ plus systemabhängige Bit-Energie) kommt ohne
  ihn aus und zeigt: er ist eine Umparametrisierung der KMS-Relation,
  keine unabhängige physikalische Eingabe (Dok.\ 325).
  Zusätzlich liefert FFGFT die Leistungskorrektur
  $P = P_{\mathrm{std}}(1-\xi\ln(M/m_P))$ = 0{,}6–1{,}4\,\%,
  für die Matzke keine Entsprechung hat.
\end{tcolorbox}

% ============================================================
\section{Vergleichstabelle}

\begin{longtable}{@{}p{3.8cm}p{4.5cm}p{4.5cm}@{}}
  \toprule
  \textbf{Größe} & \textbf{Matzke/Hyperbit} & \textbf{FFGFT} \\
  \midrule\endhead\bottomrule\endfoot
  Fundamentales Objekt & $e_i^2=+1$ \SETZ & $\xi = 4/30000$ aus Spektrum \K \\[3pt]
  Bit-Energie & $E_{\mathrm{bit}} = k_BT_P\ln 2$ (universell) \SETZ & $E_{\mathrm{bit}} = \hbar c/L$ (systemabhängig) \K \\[3pt]
  Landauer-Rolle & Physikalische Brücke: $m_{\mathrm{bit}} = m_P\ln 2/(2\pi)$ & Träger-Operation, nicht abstrakte Info (A271) \\[3pt]
  Bitwert & Universell, $\approx 0{,}11\,m_P$ & Systemabhängig, Skala $L$-abhängig \\[3pt]
  Planck-Länge & Compton $=$ Schwarzschild \B & $G = \xi^2/(4m_e)$, $\lP = \sqrt{\hbar G/c^3}$ \K \\[3pt]
  Minimal-Länge & $\lP$ & $L_0 = \xi\cdot\lP \approx \lP/7500$ \K \\[3pt]
  Stabilitätsschwelle & $n_{\mathrm{thresh}} = 6{,}41$ Bits \B & Orbifold-Fixpunktstruktur, kein Bit-Zählwert \\[3pt]
  $r_s$ ohne ART & $(2/\sqrt{\pi})\sqrt{\SBH}\,\lP$ \B & $2GM/c^2$ mit $G$ aus $\xi$ \K \\[3pt]
  $\alpha^{-1}$ & Nicht abgeleitet & $137{,}036$ aus $D_f = 3-\xi$ \K \\[3pt]
  Leptonenmassen & Nicht abgeleitet & $m_e,m_\mu,m_\tau$ aus $(r_i,p_i)$ \K \\[3pt]
  $\sin^2\theta_W$ & Nicht abgeleitet & $0{,}2308$ \K\ (Dok.\ 323) \\[3pt]
  Hawking-Mechanismus & $V/V_s$-Annihilation \B & KMS $+$ $\mathbb{Z}_3$-Selektion \K/\B\ (Dok.\ 325) \\[3pt]
  $N_c=3$ algebraisch & Furey-Identifikation \SETZ & $\mathbb{Z}_3$-Trialisierung \B\ (Dok.\ 321) \\[3pt]
  Holographie & $N$-dim $\to$ 3D \B & $T^4/\mathbb{Z}_3$-Projektion \K \\[3pt]
  Informationserhaltung & $V$-Kodierung der Voxel-Orientierung \B & Unitarität des Spektraloperators \K \\[3pt]
\end{longtable}

% ============================================================
\section{FFGFT-Einwände gegen Matzkes Landauer-Brücke}

Drei konkrete Einwände aus dem FFGFT-Korpus:

\begin{enumerate}

\item \textbf{Landauer gilt für Träger-Operationen, nicht für abstrakte Information}
  (A272, A273).
  Was Landauer 1961 beweist, ist eine untere Schranke für
  \emph{Träger-Operationen an einem thermischen Gleichgewichts-Ensemble} —
  keine universelle Aussage über abstrakte Information.
  Die Gleichsetzung $m_{\mathrm{bit}} = m_P\ln 2/(2\pi)$ behandelt das Bit als
  thermodynamisches Objekt bei Planck-Temperatur.
  Das ist eine Realisierungsannahme, kein Gesetz.

\item \textbf{Die Bit-Energie ist skalenspezifisch} (Dok.\ 302, Dok.\ 257).
  In FFGFT ist die minimale Informationseinheit das Windungsquant $\Delta w = 1$ auf $T^4$
  — skalenuniversell als \emph{Zählgröße},
  aber mit skalenspezifischer Energie $E_{\mathrm{bit}} = \hbar c/L$.
  Landauers $k_BT\ln 2$ ist der Spezialfall bei einer bestimmten Skala und Badtemperatur.
  Die Formulierung \enquote{die Bit-Energie ist $m_P\ln 2/(2\pi)$}
  impliziert eine Skala, die in FFGFT nicht ausgezeichnet ist.

\item \textbf{Die Stabilitätsschwelle ist keine rein algebraische Aussage}
  (abgeleitet aus Punkt 1 und 2).
  $n_{\mathrm{thresh}} = 2\pi/(\sqrt{2}\ln 2) = 6{,}41$ folgt aus
  Compton-Schwarzschild-Gleichung plus $m_{\mathrm{bit}} = m_P\ln 2/(2\pi)$.
  Das erste Stück (Compton = Schwarzschild) ist geometrisch \B.
  Das zweite (Landauer bei $T_P$) ist eine Setzung.
  Die scheinbare Eleganz des Ergebnisses ($\mathrm{Cl}(6)$ genau 0{,}41 Bits unter Schwelle)
  hängt an dieser Setzung.
  Eine systemabhängige Bit-Energie würde die Schwelle verschieben.

\end{enumerate}

% ============================================================
\section{Offene Brücken}

\begin{tcolorbox}[convbox, title={Punkte mit Potential für gegenseitige Befruchtung}]
  \begin{itemize}
    \item[\S] \textbf{$n_{\mathrm{thresh}}$ aus der Orbifold-Struktur}:
      Dok.\ 325 zeigt, dass Matzkes Bitwert eine Umparametrisierung der
      KMS-Relation ist. Offen bleibt die konstruktive Gegenfrage:
      Lässt sich die FFGFT-Fixpunktbedingung für Fermion-Stabilität
      in eine effektive Bit-Zähl-Aussage übersetzen — und welche
      systemspezifische Bit-Energie geht dann ein?

    \item[\S] \textbf{Casimir-Verbindung} (Dok.\ 324, R84):
      $\xi = C_2(\mathrm{SU}(3)_{\mathrm{fund}})/N_{\mathrm{Fourier}} = (4/3)/10^4$.
      Dies verbindet die algebraische Struktur von $\mathrm{Cl}(6)$ (Casimir-Invariante $C_2 = 4/3$,
      aus $N_c = 3$ \B)
      mit der geometrischen Skala der FFGFT (Torus-Modenzahl $N_F = 10^4$).
      Diese Brücke ist abgeschlossen \K\ (Dok.\ 324).

    \item[\S] \textbf{$\sin^2\theta_W$ aus $\mathrm{Cl}(6)$}:
      FFGFT leitet $\sin^2\theta_W(M_Z) = 0{,}2308$ ab \K\ (Dok.\ 323).
      Kann Matzkes $\mathrm{Cl}(6)$-Struktur einen analogen Wert erzeugen?
  \end{itemize}
\end{tcolorbox}

% ============================================================
\section{Statusbilanz}

\begin{center}
\begin{tabular}{@{}lcc@{}}
  \toprule
  \textbf{Ergebnis} & \textbf{Matzke} & \textbf{FFGFT} \\
  \midrule
  $r_s$ ohne ART & \B & \K \\
  $\lP$ aus Selbstkonsistenz & \B & \K \\
  $\mathrm{Cl}(6)$ = Fermion-Generation & \B\ (Furey) & \B\ (Orbifold) \\
  $N_c=3$ algebraisch notwendig & \SETZ & \B \\
  Bit-Energie systemabhängig & \X & \K \\
  $\alpha^{-1} = 137{,}036$ & \X & \K \\
  Leptonenmassen & \X & \K \\
  $\sin^2\theta_W$ & \X & \K \\
  Hawking-Mechanismus ausgearbeitet & \B & \K/\B\ (Dok.\ 325) \\
  $n_{\mathrm{thresh}}$ (Bit-Zahl, universell) & \B & \X\ (abgelehnt: systemabh.) \\
  Informationserhaltung & \B & \K \\
  Holographie $\neq 3D\to 2D$ & \B & \K \\
  \bottomrule
\end{tabular}
\end{center}

% ============================================================
\section*{Fazit}

Matzkes Hyperbit-Framework und FFGFT sind komplementäre Strukturen mit
bemerkenswert ähnlichen strukturellen Befunden: Fermion-Generation = algebraische
Grundeinheit, $N_c = 3$ erzwungen, Schwarzschild-Radius ohne ART.

Die tiefste Differenz liegt in der Behandlung von Bitwerten.
Matzke setzt einen universellen Bitwert über die Landauer-Brücke bei Planck-Temperatur —
eine elegante, aber nach FFGFT-Maßstäben unzulässige Universalisierung.
FFGFT stellt fest: Bitwerte sind systemabhängig ($E_{\mathrm{bit}} = \hbar c/L$),
Landauer gilt für Träger-Operationen eines thermischen Ensembles,
und die Stabilitätsschwelle $n_{\mathrm{thresh}}$ wäre dann keine universelle Zahl.

Dies ist kein Einwand gegen die algebraische Schönheit von Matzkes Ansatz,
sondern ein theoretischer Unterschied in der Grundlegung.
Dok.\ 325 macht ihn konstruktiv: die vollständige Hawking-Physik —
Temperatur, Selektion, Träger, Informationserhalt — folgt aus
$\tilde{T}\cdot m = 1$, der $\mathbb{Z}_3$-Struktur und der
systemabhängigen Bit-Energie, ohne universellen Bitwert.
Matzkes $n_{\mathrm{thresh}} = 6{,}41$ hängt damit nachweislich an der
Bitwert-Parametrisierung, nicht an der Temperaturphysik.

\begin{tcolorbox}[ffgftbox, title={Hauptbefund}]
  Matzke und FFGFT konvergieren auf \textbf{Fermion-Generation = fundamentale algebraische Einheit}.
  Sie divergieren bei \textbf{Bitwerten}: universell (Matzke, Landauer bei $T_P$)
  vs.\ systemabhängig (FFGFT, $E_{\mathrm{bit}} = \hbar c/L$).
  Dieser Unterschied ist nicht kosmetisch — er betrifft die Stabilitätsschwelle
  und die Masse eines Bits. Die Hawking-Temperaturformel selbst ist davon
  unberührt: beide Routen geben dieselbe Zahl, aber nur die FFGFT-Route
  (Dok.\ 325) kommt ohne universellen Bitwert aus.
\end{tcolorbox}

% ============================================================
\section*{Bezug zum FFGFT-Korpus}

\begin{center}\small
\begin{tabular}{@{}lll@{}}
  \toprule
  \textbf{Thema} & \textbf{Status} & \textbf{Dok.} \\
  \midrule
  Landauer: Phasenraum-Buchhaltung & \K & A271 \\
  Träger vs.\ Information (Landauer-Kritik) & \K & A272 \\
  Bit-Energie systemabhängig (Rechenkugel) & \K & A273 \\
  Informationseinheit $\Delta w = 1$, $E_{\mathrm{bit}} = \hbar c/L$ & \K & Dok.\ 257 \\
  Elementarzellen, Bit-Kapazität, Partitionsinvarianz & \K & Dok.\ 302 \\
  $G = \xi^2/(4m_e)$, $\lP$, $L_0$ & \K & Dok.\ 180 \\
  $\mathrm{SU}(3)_c$ aus $\mathbb{Z}_3$-Trialisierung & \B & Dok.\ 321 \\
  $\sin^2\theta_W(M_Z) = 0{,}2308$ & \K & Dok.\ 323 \\
  $\xi = C_2(\mathrm{SU}(3))/N_{\mathrm{Fourier}}$ & \K & Dok.\ 324, R84 \\
  Hawking-Mechanismus (KMS, $\mathbb{Z}_3$, $n_4{=}0$) & \K/\B & Dok.\ 325 \\
  Spektraltheorie, Hilbertraum-Einbettung & \K & Dok.\ 322 \\
  Leptonenmassen aus $(r_i, p_i)$ & \K & Dok.\ 006, 046, 320 \\
  $E = mc^2 = E = m$ (Einheitenidentität) & \K & Dok.\ 077 \\
  \bottomrule
\end{tabular}
\end{center}

\section*{Literatur}

\begin{enumerate}
  \item D.\ Matzke, \textit{Black Holes from Hyperbits}, IPI Annual Conference, August 2026.
  \item J.\ Pascher, A271 \textit{Landauer und die Phasenraum-Buchhaltung}, FFGFT-Korpus, 2026.
  \item J.\ Pascher, A272 \textit{Träger und Information}, FFGFT-Korpus, 2026.
  \item J.\ Pascher, A273 \textit{Die Rechenkugel}, FFGFT-Korpus, 2026.
  \item J.\ Pascher, Dok.\ 257 \textit{Informationseinheit auf $T^4$}, FFGFT-Korpus.
  \item J.\ Pascher, Dok.\ 302 \textit{Elementarzellen und Bit-Kapazität}, FFGFT-Korpus, 2026.
  \item J.\ Pascher, Dok.\ 180 \textit{Derivation of $L_0$}, FFGFT-Korpus, 2026.
  \item J.\ Pascher, Dok.\ 321/322/323/324, FFGFT-Korpus, 2026.
  \item J.\ Pascher, Dok.\ 325 \textit{Der FFGFT-Hawking-Mechanismus}, FFGFT-Korpus, 2026.
  \item R.\ Landauer, \textit{Irreversibility and heat generation}, IBM J.\ Res.\ Dev.\ \textbf{5}, 183, 1961.
  \item C.\ Furey, \textit{Standard Model Physics from an Algebra?}, Diss., Waterloo, 2015.
  \item J.\ Bekenstein, Phys.\ Rev.\ D \textbf{7}, 2333, 1973.
\end{enumerate}

\vspace{6pt}\noindent\small
\textcopyright\ 2026 Johann Pascher · \href{https://creativecommons.org/licenses/by/4.0/}{CC BY 4.0}
